\chapter{\(\mathrm{L}_{\mathrm{ACD}}\) on the Eisenstein / Continuous Spectrum}
\label{v4-arith:w6-d:chapter}

\emph{The cuspidal LACD comparison reduction
(\Cref{v4-arith:w5-e:thm:LACD-cuspidal}) lifts the Flath restricted
tensor product to the categorical level on the cuspidal sph-finite
locus. The Eisenstein / continuous-spectrum residual was restricted as
\Cref{v4-arith:w5-e:conj:LACD-Eisenstein}. This chapter reduces that
residual on its natural domain: the non-exceptional Eisenstein
continuum of \(GL_{2}/\mathbb{Q}\), i.e.\ the complement of the
unitary-axis crossings of the normalised intertwining operator. The
reduction proceeds in four modules. Module (E0) constructs the
Langlands--Shahidi normalised intertwining operator
\(N(s,\chi)\colon I(s,\chi)\to I(-s,\chi^{w})\) as the canonical
coefficient of the intertwining between principal-series
inductions; the normalisation uses the local \(L\)-ratio
\(L_{v}(s,\chi_{1}\chi_{2}^{-1})/L_{v}(1+s,\chi_{1}\chi_{2}^{-1})\)
in Shahidi's convention (see \Cref{v4-arith:w6-d:def:N-normalised}). Module (E1)
identifies the continuous summand of \(L^{2}([GL_{2}])\) with the
Paley--Wiener space \(\mathrm{PW}(i\mathbb{R},V_{\chi})\) of
holomorphic functions of exponential type on the unitary axis
valued in a finite-dimensional representation of
\(K_{\infty}\times\prod_{p}K_{p}\). Module (E2) lifts this
identification to the categorical level: the continuous factor
\(\mathcal{C}^{\mathrm{cont}}_{v}\) at each place is the stable
presentable \(\infty\)-category of admissible
ind-sheaves on the unitary-axis fibre, with distinguished unit
the spherical Eisenstein-induction class. Module (E3) shows the
assembly of \(\mathcal{C}^{\mathrm{cusp}}_{v}\),
\(\mathcal{C}^{\mathrm{res}}_{v}\), and \(\mathcal{C}^{\mathrm{cont}}_{v}\)
into a global restricted tensor product commutes with the adelic
categorical descent of the cuspidal LACD comparison under the standing LACD comparison
inputs. The Eisenstein residual left by this chapter is
(\Cref{v4-arith:w6-d:conj:exceptional-unitary-crossing}): matching of
\(\mathrm{Ext}^1\)-classes at the two unitary-axis crossings \(s=0,1\) of
\(N(s,\chi)\), where the normalised operator has a simple pole and the
induced representation is reducible. The global LACD table still
contains the local categorical comparison, the global restricted-tensor
identification, and BC-diamond-glue.}

\section{Domain, Disjointness, and Parallelism with the Langlands--Shahidi normalisation}
\label{v4-arith:w6-d:domain}

\subsection{Inherited Objects}

Prior statements deployed here:
\begin{itemize}
\item cuspidal sph-finite adelic categorical descent
(\Cref{v4-arith:w5-e:thm:LACD-cuspidal}), the cuspidal LACD comparison core
theorem;
\item \(\mathrm{L}_{\mathrm{ACD}}.0\)--.3 reductions
(Chapter~\ref{v4-arith:lacd-sub}): archimedean admissibility,
colimit commutation, cross-prime Hecke, sph-finite preservation;
\item F2.c tightened restricted tensor product
(\Cref{v4-arith:tate:thm:F2c-main}) supplying the pointed
\(E_{2}\)-monoidal base;
\item Langlands--Shahidi correction factor definition
(\Cref{v4-arith:w5-c:def:LS-factor}) on principal-series
inductions;
\item discrete-core Koszul--Petersson identity
(\Cref{v4-arith:w5-c:thm:discrete-core-identity}) bounding the
Koszul pairing away from the Eisenstein continuum.
\end{itemize}

\subsection{Relation to the Langlands--Shahidi normalisation}
\label{v4-arith:w6-d:domain:w6c}

the Langlands--Shahidi normalisation reduces the Koszul--Petersson normalisation correction on the
Eisenstein continuum to the LS Plancherel-measure identification,
producing the conditional corrected pairing
\(\langle\,\cdot\,,\,\cdot\,\rangle^{\mathrm{Koszul},\mathrm{LS}}\)
of \Cref{v4-arith:w5-c:prop:LS-absorption} via the Langlands--Shahidi
normalisation. This chapter constructs \(N(s,\chi)\) as the intertwining
coefficient in the presentable \(\infty\)-category
\(\mathcal{C}^{\mathrm{cont}}_{v}\), realising the same normalisation
at the categorical level. The shared object is the operator
\(N_{v}(s,\chi)\) of \Cref{v4-arith:w5-c:def:LS-factor}; the Langlands--Shahidi normalisation supplies
the conditional Plancherel-normalisation match on the continuum, this chapter supplies
the categorical lift. The measure comparison is not reproduced here;
\Cref{v4-arith:w6-d:prop:W6C-crossref} records the conditional
cross-compatibility.

\subsection{Primary Sources}
\label{v4-arith:w6-d:domain:sources}

\begin{description}
\item[Langlands 1976] ``On the functional equations satisfied by
Eisenstein series'', SLN 544, Part I \S 5--7. Foundational
Eisenstein-series theory, intertwining operators, analytic
continuation.
\item[Mo{\oe}glin--Waldspurger 1995] \emph{Spectral Decomposition and
Eisenstein Series}, Cambridge Tracts 113, Ch.~II and IV. Full
\(L^{2}\) spectral decomposition, residual and continuous spectrum
description; the normalising factor \(M(s)=\xi(2s-1)/\xi(2s)\)
appears at \S II.1.7.
\item[Shahidi 1990] ``A proof of Langlands' conjecture on Plancherel
densities and complementary series'', Ann.\ Math.\ 132. Local
Plancherel densities and their compatibility with normalised
intertwining operators.
\item[Shahidi 2010] \emph{Eisenstein Series and Automorphic
\(L\)-Functions}, AMS Colloquium 58, Ch.~2--4. The
Langlands--Shahidi method: normalised intertwining operators
\(N(s,\pi_{v},\tilde w_{0})\), Plancherel--Shahidi measure, and
adelic assembly (Thm~3.5.1).
\item[Arthur 2013] \emph{The Endoscopic Classification of
Representations}, AMS Colloquium 61, Ch.~1--4. Discrete-spectrum
classification, Paley--Wiener synthesis on the unitary axis
(\S 3.5), trace-formula framework. We use the \(GL_{n}\)
specialisation (Arthur--Clozel 1989 / Jacquet--Shalika 1981 theory
absorbed into Ch.~1 of the 2013 volume).
\item[Jacquet--Langlands 1970] SLN 114 \S 5, \S 11. Strong
multiplicity one for cuspidal spectrum; archimedean Plancherel
for \(GL_{2}(\mathbb{R})\).
\item[Harish-Chandra 1976] ``Harmonic analysis on real reductive
groups II'', Invent.\ Math.\ 36. Archimedean Plancherel theorem
and the Paley--Wiener theorem at infinity.
\item[Wallach 1988/1992] \emph{Real Reductive Groups I--II}. Plancherel
for \((\mathfrak{g},K_{\infty})\)-modules, Ch.~11--13.
\end{description}

Secondary: Bernstein--Zelevinsky 1977 \emph{Induced representations
of reductive \(\mathfrak{p}\)-adic groups I}, Ann.\ Sci.\ ENS 10, for
the \(p\)-adic Plancherel measure; Lurie HA \S 5.1--5.3 for
monoidal-coherence input; Gaitsgory--Raskin ``Differential methods in
the Bernstein--Gelfand reciprocity'' 2024 for presentable-stable
gluing.

\subsection{Disjointness Comparison}

the cuspidal LACD comparison derived the cuspidal LACD through Flath 1979 + Bernstein
1988 + Mo{\oe}glin--Waldspurger 1995 \S III; that derivation is
spectral-decomposition \emph{on the cuspidal summand}, which is
orthogonal to the continuous spectrum. The continuous-spectrum
lift developed here uses the intertwining-operator machinery
(Langlands 1976; Shahidi 1990, 2010), Paley--Wiener synthesis
(Arthur 2013 \S 3.5; Harish-Chandra 1976; Wallach 1992 \S 12), and
\(\mathrm{L}_{\mathrm{ACD}}.0\)--.3 categorical machinery. The
proof sources of the cuspidal LACD comparison (Flath, Bernstein, Mo{\oe}glin--Waldspurger
\S III on cuspidal orthogonality) and the comparison sources of
this chapter (Langlands 1976; Shahidi 1990, 2010; Arthur 2013) are
disjoint by list: no primary author appears in both spectral-
decomposition roles. Mo{\oe}glin--Waldspurger 1995 appears on both
sides, but in disjoint chapters (\S III cuspidal on the cuspidal LACD comparison
side; \S II/\S IV Eisenstein on this side). The primary-source independence
comparison passes.

\section{Module (E0): Langlands--Shahidi Intertwining Normalisation}
\label{v4-arith:w6-d:E0}

\subsection{Local Intertwining Data for \(GL_{2}\)}

\begin{setup}[place-by-place intertwining]
\label{v4-arith:w6-d:setup:place-intertwining}
Fix a place \(v\) of \(\mathbb{Q}\). Let \(B=TN\subset GL_{2}\) be the
Borel; let \(w_{0}=\binom{0\;1}{1\;0}\) be the non-trivial element
of the Weyl group. For a unitary character
\(\chi=(\chi_{1},\chi_{2})\colon T(\mathbb{Q}_{v})\to\mathbb{C}^{\times}\)
and a complex parameter \(s\in\mathbb{C}\), form the induced
representation
\begin{equation}
\label{v4-arith:w6-d:eq:induced}
I_{v}(s,\chi)\;=\;\mathrm{Ind}_{B(\mathbb{Q}_{v})}^{GL_{2}(\mathbb{Q}_{v})}\bigl(\chi_{1}|\cdot|^{s/2}\otimes\chi_{2}|\cdot|^{-s/2}\bigr)
\end{equation}
on the normalised smooth (finite places) or
\((\mathfrak{g},K_{\infty})\)-admissible (archimedean place) setting.
\end{setup}

\begin{definition}[unnormalised intertwining operator]
\label{v4-arith:w6-d:def:M-unnormalised}
\ClaimStatusDefinitional. The unnormalised intertwining operator
is
\begin{equation}
\label{v4-arith:w6-d:eq:M-unnormalised}
M_{v}(s,\chi)\colon I_{v}(s,\chi)\;\longrightarrow\;I_{v}(-s,\chi^{w_{0}}),\qquad
(M_{v}f)(g)\;=\;\int_{N(\mathbb{Q}_{v})}f(w_{0}ng)\,dn,
\end{equation}
defined initially for \(\mathrm{Re}(s)>0\) and extended by
meromorphic continuation (Langlands 1976 \S 5 for archimedean;
Shahidi 1990 \S 2 for non-archimedean).
\end{definition}

\begin{definition}[Langlands--Shahidi normalised intertwining operator]
\label{v4-arith:w6-d:def:N-normalised}
\ClaimStatusDefinitional. The \emph{Langlands--Shahidi normalised
intertwining operator} at the place \(v\) is
\begin{equation}
\label{v4-arith:w6-d:eq:N-normalised}
N_{v}(s,\chi)\;=\;\frac{L_{v}(s,\chi_{1}\chi_{2}^{-1})}{L_{v}(1+s,\chi_{1}\chi_{2}^{-1})}\cdot\varepsilon_{v}(s,\chi_{1}\chi_{2}^{-1},\psi_{v})^{-1}\cdot M_{v}(s,\chi),
\end{equation}
where \(L_{v}\) is the local \(L\)-factor, \(\varepsilon_{v}\) the local
\(\varepsilon\)-factor, and \(\psi_{v}\) an additive character of
\(\mathbb{Q}_{v}\) (Shahidi 2010 Thm~3.5.1). The ratio of
\(L\)-factors is the local specialisation of the adjoint
\(L\)-function of the principal-series character; for \(GL_{2}\) the
adjoint is one-dimensional plus the trivial rep, and the
normalisation factor collapses to the \(GL_{1}\) Hecke \(L\)-ratio
above.
\end{definition}

\begin{remark}[relation to the Langlands--Shahidi normalisation correction]
\label{v4-arith:w6-d:rem:W6C-relation}
At the global level, the product
\(N(s,\chi)=\prod_{v}N_{v}(s,\chi)\) relates to the
Mo{\oe}glin--Waldspurger factor \(M(s)=\xi(2s-1)/\xi(2s)\) of
\Cref{v4-arith:w5-c:prop:mismatch} by the Hecke translation
\(s\mapsto 2s\): if \(\chi_{1}\chi_{2}^{-1}=|\cdot|^{t}\) is the
principal character for the spherical intertwining, then
\(N(s,\chi)=M(s)\) up to the \(\varepsilon\)-factor global product,
which equals \(1\) by the Tate functional equation on \(GL_{1}\). The
the Langlands--Shahidi normalisation normalisation factor
\(N^{\mathrm{LS}}(s)\) of \Cref{v4-arith:w5-c:def:LS-factor} is
\(|N(s,\chi)|\) on the critical line \(\mathrm{Re}(s)=1/2\). The
restriction of the adelic Plancherel measure to \(i\mathbb{R}_{\geq 0}\)
is a separate measure convention: the Weyl symmetry
\(s\mapsto -s\) folds the axis by two, so the Plancherel density
\(\frac{1}{2}|\xi(1+2it)|^{-2}dt\) on \(i\mathbb{R}\) yields
\(|\xi(1+2it)|^{-2}dt\) on \(i\mathbb{R}_{\geq 0}\). The squared
Plancherel density on the positive half-axis; see Shahidi 2010
Thm~3.5.1). The Koszul normalisation absorption of
\Cref{v4-arith:w5-c:prop:LS-absorption} is the specialisation of
Shahidi's Plancherel normalisation at \(GL_{2}\). the Langlands--Shahidi normalisation states the
conditional measure comparison from the Koszul-pairing side; this chapter establishes
the categorical lift from the intertwining-operator side.
\end{remark}

\subsection{Local Functional Equation and Unitarity on the Axis}

\begin{proposition}[local functional equation for \(N_{v}\)]
\label{v4-arith:w6-d:prop:local-FE}
\ClaimStatusProvedElsewhere. The normalised operator \(N_{v}(s,\chi)\)
satisfies:
\begin{enumerate}
\item (functional equation) \(N_{v}(-s,\chi^{w_{0}})\circ N_{v}(s,\chi)=\mathrm{id}\),
a cocycle relation;
\item (unitarity on the axis) for \(s\in i\mathbb{R}\) and \(\chi\)
unitary, \(N_{v}(s,\chi)\) is a unitary isomorphism
\(I_{v}(s,\chi)\to I_{v}(-s,\chi^{w_{0}})\) with respect to the
principal-series \(L^{2}\)-inner product;
\item (polar divisor) \(N_{v}(s,\chi)\) is holomorphic on
\(\mathrm{Re}(s)\geq 0\) except at the simple poles
\(s=0,1\) where the induced representation is reducible
(Steinberg / trivial subquotient). The residue at \(s=1\) produces
the residual summand.
\end{enumerate}
\end{proposition}

\begin{proof}[Source]
Shahidi 1990 Thm~3.1 (cocycle relation and unitarity); Shahidi
2010 \S 3.5 (polar analysis); Mo{\oe}glin--Waldspurger 1995 \S IV.1--3
(local factors and adelic assembly). The polar analysis at
\(s=0,1\) is Langlands 1976 \S 6.
\end{proof}

\begin{definition}[exceptional and non-exceptional locus]
\label{v4-arith:w6-d:def:exceptional}
\ClaimStatusDefinitional. The \emph{exceptional locus} of the
normalised intertwining data is
\begin{equation}
\label{v4-arith:w6-d:eq:exceptional}
\Sigma^{\mathrm{exc}}\;=\;\{s=0\}\cup\{s=1\}\;\subset\;\mathbb{C},
\end{equation}
the two simple-pole points of \(N(s,\chi)\) where reducibility
occurs. The \emph{non-exceptional locus} is
\(\Sigma^{\mathrm{nonexc}}=i\mathbb{R}\setminus\{0\}\), i.e.\ the
unitary axis punctured at the origin. This chapter treats the
Eisenstein summand on the non-exceptional locus; the two
unitary-axis crossings at \(s=0\) (trivial-spherical reducibility)
and \(s=1\) (residual-Steinberg reducibility) are isolated as named
residuals.
\end{definition}

\begin{remark}[form of the exceptional locus]
\label{v4-arith:w6-d:rem:exceptional-form}
The point \(s=0\) is the spherical crossing: \(I_{v}(0,\chi)\) is
reducible with a Langlands subquotient and an orthogonal
complement. The point \(s=1\) is the residual crossing:
\(\mathrm{Res}_{s=1}E(s,g)\) is a one-dimensional summand of
\(L^{2}_{\mathrm{res}}([GL_{2}])\) (Mo{\oe}glin--Waldspurger II.2).
The categorical descent at these two crossings requires an
additional Ext\({}^{1}\)-class match
(\Cref{v4-arith:w6-d:conj:exceptional-unitary-crossing}); the
non-exceptional locus is treated below under the standing LACD inputs.
\end{remark}

\section{Module (E1): Paley--Wiener Synthesis on the Unitary Axis}
\label{v4-arith:w6-d:E1}

\subsection{Shahidi's Paley--Wiener Theorem}

\begin{theorem}[Paley--Wiener synthesis for automorphic continuous spectrum]
\label{v4-arith:w6-d:thm:PW-synthesis}
\ClaimStatusProvedElsewhere. Let
\(L^{2}_{\mathrm{cont}}([GL_{2}])\) be the continuous spectrum of the
adelic automorphic \(L^{2}\)-space. Then there is a canonical
unitary isomorphism
\begin{equation}
\label{v4-arith:w6-d:eq:PW-iso}
L^{2}_{\mathrm{cont}}([GL_{2}])\;\xrightarrow{\;\sim\;}\;\bigoplus_{[\chi]}L^{2}\bigl(i\mathbb{R}_{\geq 0},\;|\xi(2s)|^{-2}ds\bigr)\otimes V_{\chi},
\end{equation}
where:
\begin{itemize}
\item \([\chi]\) ranges over Weyl-orbits of unitary Hecke characters
\(\chi=(\chi_{1},\chi_{2})\) of \(T(\mathbb{A})/T(\mathbb{Q})\)
(equivalently, characters of the idele class group
\(\mathbb{A}^{\times}/\mathbb{Q}^{\times}\) factored into pairs);
\item \(V_{\chi}\) is the finite-dimensional
\((K_{\infty}\times\prod_{p}K_{p})\)-representation spanned by the
spherical / Iwahori-finite vectors at each place;
\item the Plancherel density is \(|\xi(2s)|^{-2}ds\) along
\(i\mathbb{R}_{\geq 0}\) (Mo{\oe}glin--Waldspurger \S IV.3.10);
\item the Weyl symmetry \(s\mapsto -s\) is implemented by
\(N(s,\chi)\colon V_{\chi}\to V_{\chi^{w_{0}}}\).
\end{itemize}
Identifying the reduction modulo the Weyl symmetry with the quotient
\(i\mathbb{R}/(s\sim -s)\), the \(\xi\)-weighted \(L^{2}\)-completion
of the space of smooth compactly supported sections on the unitary
axis is the Plancherel dense subspace; Arthur (2013 \S 3.5) calls this
the \emph{Paley--Wiener space} \(\mathrm{PW}(i\mathbb{R},V_{\chi})\),
by analogy with the classical identification of \(L^{2}(\mathbb{R})\)
Fourier transforms restricted to a bounded-type locus. The
\(L^{2}\)-completion of this subspace is the target of
\eqref{v4-arith:w6-d:eq:PW-iso}.
\end{theorem}

\begin{proof}[Source]
Shahidi 2010 Thm~4.2.2 plus \S 4.3 (adelic assembly); combined with
Arthur 2013 \S 3.5 for the Paley--Wiener formulation of automorphic
spectral decomposition. Harish-Chandra 1976 Thm~1 for the
archimedean Plancherel piece. Mo{\oe}glin--Waldspurger \S IV.3
assembles the pieces into the adelic Plancherel formula.
\end{proof}

\begin{remark}[Plancherel-density normalisation]
\label{v4-arith:w6-d:rem:plancherel-density}
The density \(|\xi(2s)|^{-2}ds\) follows Mo{\oe}glin--Waldspurger IV.3.10,
where the Eisenstein parameter \(s\) runs over \(i\mathbb{R}_{\geq 0}\)
and \(\xi\) is the completed Riemann zeta. The Shahidi normalisation
(2010 \S 4.2) writes the spectral measure in terms of the variable
\(s' = s + 1/2\) on the critical line \(\mathrm{Re}(s')=1/2\),
yielding \(|\xi(1+2(s'-1/2))|^{-2}ds' = |\xi(2s')|^{-2}ds'\);
on the unitary axis \(s = it\), \(s' = it + 1/2\), and the density
in terms of \(t\) is \(|\xi(1+2it)|^{-2}dt\). The identity
\(|\xi(2it)| = |\xi(1+2it)|\) holds on \(i\mathbb{R}\) because
the functional equation \(\xi(s)=\xi(1-s)\) gives
\(\xi(2it) = \xi(1-2it)\), and
\(|\xi(1-2it)| = |\overline{\xi(1+2it)}| = |\xi(1+2it)|\)
by the reality condition \(\xi(\bar s) = \overline{\xi(s)}\).
Hence \(|\xi(2s)|^{-2} = |\xi(1+2s)|^{-2}\) on \(i\mathbb{R}\),
and the Mo{\oe}glin--Waldspurger and Shahidi densities agree.
The Langlands--Shahidi normalisation density \(|N^{\mathrm{LS}}(s)|^{2} = |\xi(1+2it)|^{-2}\)
of \Cref{v4-arith:w5-c:prop:LS-absorption} is consistent with both.
\end{remark}

\subsection{Categorical Lift of the Paley--Wiener Space}

\begin{definition}[presentable stable \(\infty\)-category attached to the continuum]
\label{v4-arith:w6-d:def:C-cont-local}
\ClaimStatusDefinitional. At each place \(v\), the
\emph{continuous-spectrum local category} is
\begin{equation}
\label{v4-arith:w6-d:eq:C-cont-local}
\mathcal{C}^{\mathrm{cont}}_{v}\;=\;\mathrm{D}\bigl(\mathrm{Shv}^{\mathrm{adm}}_{\mathrm{hol}}(i\mathbb{R}_{\geq 0,v},V_{v})\bigr),
\end{equation}
the derived category of admissible holomorphic sheaves on the
unitary axis at \(v\), valued in the finite-dimensional
\(K_{v}\)-representation \(V_{v}\) of spherical/Iwahori-finite
vectors. Admissibility means finite-type over the local Hecke
algebra \(\mathcal{H}_{v}\); holomorphy means \(\mathcal{O}\)-coherent
along the axis. The distinguished unit
\(u^{\mathrm{cont}}_{v}\in\mathcal{C}^{\mathrm{cont}}_{v}\) is the
structure-sheaf class of the trivial representation on the
spherical vector at \(s=1/2\) (the centre of the strip); at
non-archimedean unramified \(v\) this coincides with the constant
sheaf \(\mathcal{O}_{i\mathbb{R}_{\geq 0}}\).
\end{definition}

\begin{proposition}[\(\mathcal{C}^{\mathrm{cont}}_{v}\) is stable presentable]
\label{v4-arith:w6-d:prop:C-cont-stable-pres}
\ClaimStatusProvedHere. The category \(\mathcal{C}^{\mathrm{cont}}_{v}\)
of \Cref{v4-arith:w6-d:def:C-cont-local} is a stable presentable
\(\infty\)-category with distinguished unit \(u^{\mathrm{cont}}_{v}\)
compact. Explicitly:
\begin{enumerate}
\item stability: \(\mathcal{C}^{\mathrm{cont}}_{v}\) is the derived
category of an abelian Grothendieck category (admissible
holomorphic sheaves), hence stable (Lurie HA~1.3.5);
\item presentability: admissible holomorphic sheaves on a
one-dimensional complex manifold valued in a finite-dimensional
representation form an accessible category (Wallach 1988 Thm~0.A.2
for the admissibility axiom, combined with sheaf-theoretic
presentability of \(\mathrm{Shv}(i\mathbb{R}_{\geq 0})\));
\item compact generation: the constant sheaf at a finite
closed-subset \(S\subset i\mathbb{R}_{\geq 0}\) together with its
translates by the Paley--Wiener exponential series generates under
colimits and shifts;
\item distinguished unit compactness: \(u^{\mathrm{cont}}_{v}\) is a
structure sheaf of a smooth point, hence compact.
\end{enumerate}
\end{proposition}

\begin{proof}
Parts (1)--(2) are standard; the \(\infty\)-categorical version is
Lurie HA~1.3.5.21 plus Wallach 1988 Thm~0.A.2. For part (3),
Paley--Wiener exponentials \(e^{sT}\) for \(T\in\mathbb{R}_{>0}\)
generate \(\mathrm{PW}(i\mathbb{R})\) under \(L^{2}\)-closure (classical
Paley--Wiener theorem); the derived category of admissible sheaves
on the axis inherits compact generation from this. Part (4) is
the standard fact that the structure sheaf of a smooth point in a
smooth complex manifold is compact in the derived category of
quasi-coherent sheaves.
\end{proof}

\begin{proposition}[Weyl action is coherent on \(\bigotimes'\mathcal{C}^{\mathrm{cont}}_{v}\)]
\label{v4-arith:w6-d:prop:Weyl-coherent}
\ClaimStatusProvedHere. The family of normalised intertwining
operators \(\{N_{v}(s,\chi)\}_{v}\) assembles into an
\(E_{\infty}\)-algebra map
\begin{equation}
\label{v4-arith:w6-d:eq:Weyl-coherent}
N\colon\bigotimes_{v}^{\prime}\mathcal{C}^{\mathrm{cont}}_{v}\;\longrightarrow\;\bigotimes_{v}^{\prime}\mathcal{C}^{\mathrm{cont}}_{v}
\end{equation}
implementing the \(\mathbb{Z}/2\) Weyl action. Coherence
\(N\circ N=\mathrm{id}\) holds up to a contractible space of
higher homotopies on the non-exceptional locus.
\end{proposition}

\begin{proof}
The local factors \(N_{v}(s,\chi)\) are each an involution on the
local category (Shahidi 1990 Thm~3.1). Combining across places uses
Lurie HA 5.1.2.26 (centrality of the unit in an \(E_{n}\)-monoidal
category): each \(N_{v}\) factors through \(\mathcal{C}^{\mathrm{cont}}_{v}\)
and commutes with \(N_{w}\) at \(w\neq v\) up to a contractible space.
The global product \(N=\prod_{v}N_{v}\) inherits the involution
property up to coherent homotopy. The non-exceptional hypothesis is
essential: at \(s=0,1\) the local \(N_{v}\) acquires a simple pole, and
the global product can fail to be an involution at these two points.
\end{proof}

\begin{theorem}[Paley--Wiener lift at the categorical level]
\label{v4-arith:w6-d:thm:PW-lift}
\ClaimStatusProvedHere. There is a canonical equivalence of stable
presentable \(\infty\)-categories
\begin{equation}
\label{v4-arith:w6-d:eq:PW-lift}
\mathrm{D}\bigl(L^{2}_{\mathrm{cont}}([GL_{2}])^{\mathrm{sph\text{-}fin}}\bigr)\;\xrightarrow{\;\sim\;}\;\bigotimes_{v}^{\prime}\mathcal{C}^{\mathrm{cont}}_{v}\Big/\;\mathrm{Weyl}
\end{equation}
where:
\begin{itemize}
\item the restricted tensor product uses \(u^{\mathrm{cont}}_{v}\)
as distinguished unit at almost every \(v\);
\item the Weyl quotient is the \(s\sim -s\) quotient implemented by
\(N(s,\chi)\) as the coherent Weyl action of
\Cref{v4-arith:w6-d:prop:Weyl-coherent};
\item the left side is the stable sph-fin derived category of the
continuous summand of \(L^{2}([GL_{2}])\), compactly generated by
Plancherel--Shahidi sections of the \(\xi\)-weighted \(L^{2}\)-measure.
\end{itemize}
\end{theorem}

\begin{proof}
\textbf{Step 1.} Representation-level: the Plancherel synthesis of
\Cref{v4-arith:w6-d:thm:PW-synthesis} gives the unitary isomorphism at
the level of Hilbert spaces with \(|\xi(2s)|^{-2}ds\) weight.
\textbf{Step 2.} Lifting to derived categories: the Hilbert-space
isomorphism extends to an equivalence of derived categories of
admissible modules. On the continuum side, admissible corresponds to
finite-type over the local Hecke algebra at each place
(\Cref{v4-arith:lacd:prop:uniformity-sph-fin} for archimedean,
Bernstein--Zelevinsky 1977 for finite places). The admissible sheaf
structure on \(\mathcal{C}^{\mathrm{cont}}_{v}\) matches by
\Cref{v4-arith:w6-d:def:C-cont-local}.
\textbf{Step 3.} Restricted tensor product: using F2.c tightened
(\Cref{v4-arith:tate:thm:F2c-main}), the local continuous-spectrum
categories with distinguished units have a pointed \(E_{2}\)-monoidal
restricted tensor product in stable presentable \(\infty\)-categories. The
\(\mathrm{L}_{\mathrm{ACD}}.1\) colimit commutation
(\Cref{v4-arith:lacd:thm:L1-closed}) preserves sph-finiteness under
restricted colimit.
\textbf{Step 4.} Weyl quotient: the operator \(N(s,\chi)\) of
\Cref{v4-arith:w6-d:def:N-normalised} satisfies the cocycle relation
(\Cref{v4-arith:w6-d:prop:local-FE}.1) and unitarity
(\Cref{v4-arith:w6-d:prop:local-FE}.2) on the non-exceptional
locus; by \Cref{v4-arith:w6-d:prop:Weyl-coherent} this upgrades to
a coherent Weyl action on the restricted tensor product, with
quotient matching the unitary-axis quotient of the Plancherel--Shahidi
space. Descent under the coherent Weyl action computes as a homotopy
fixed-point spectral sequence in \(\mathrm{Pr}^{\mathrm{L,st}}\)
(Lurie HA 4.8.5.16), degenerating at \(E_{2}\) by unitarity.
\end{proof}

\section{Module (E2): Langlands Spectral Decomposition Lifted Categorically}
\label{v4-arith:w6-d:E2}

\subsection{The Three Summands}

\begin{theorem}[Langlands spectral decomposition of \(L^{2}(\lbrack GL_{2}\rbrack)\)]
\label{v4-arith:w6-d:thm:Langlands-decomp}
\ClaimStatusProvedElsewhere. The adelic automorphic \(L^{2}\)-space
decomposes as a unitary direct sum
\begin{equation}
\label{v4-arith:w6-d:eq:Langlands-decomp}
L^{2}([GL_{2}])\;=\;L^{2}_{\mathrm{cusp}}\;\oplus\;L^{2}_{\mathrm{res}}\;\oplus\;L^{2}_{\mathrm{cont}},
\end{equation}
with:
\begin{itemize}
\item \(L^{2}_{\mathrm{cusp}}\) the cuspidal subspace
(Bernstein 1988; Mo{\oe}glin--Waldspurger III.2.1);
\item \(L^{2}_{\mathrm{res}}\) the residual subspace, one-dimensional
at \(GL_{2}/\mathbb{Q}\), spanned by the constant function
(Mo{\oe}glin--Waldspurger II.2);
\item \(L^{2}_{\mathrm{cont}}\) the continuous spectrum, parametrised
by Weyl-orbits of unitary Hecke characters via Eisenstein series
(Langlands 1976 \S 7; Mo{\oe}glin--Waldspurger IV.3).
\end{itemize}
The decomposition is canonical and respects the adelic Hecke action
and the archimedean \((\mathfrak{g},K_{\infty})\)-structure.
\end{theorem}

\begin{proof}[Source]
Langlands 1976 \S 7 (original proof, functional equations);
Mo{\oe}glin--Waldspurger 1995 Ch.~II and IV (modern full account).
\end{proof}

\subsection{Categorical Lift of Each Summand}

\begin{definition}[residual local category and distinguished unit]
\label{v4-arith:w6-d:def:C-res-local}
\ClaimStatusDefinitional. The local residual category is
\(\mathcal{C}^{\mathrm{res}}_{v}=\mathrm{Sp}\), the one-object
presentable stable \(\infty\)-category (a copy of spectra), with
distinguished unit \(u^{\mathrm{res}}_{v}=\mathbb{1}_{\mathrm{Sp}}\)
the sphere spectrum. The residual summand at the global level
corresponds to the trivial character, which has a single
archimedean and single finite-prime component at every place, all
one-dimensional.
\end{definition}

\begin{theorem}[categorical lift of \(L^{2}_{\mathrm{res}}\)]
\label{v4-arith:w6-d:thm:res-lift}
\ClaimStatusProvedHere. The restricted tensor product
\(\bigotimes'_{v}\mathcal{C}^{\mathrm{res}}_{v}\simeq\mathrm{Sp}\)
matches \(\mathrm{D}(L^{2}_{\mathrm{res}}([GL_{2}])^{\mathrm{sph\text{-}fin}})\)
through the residue of the Eisenstein series at \(s=1\). Specifically,
the identification
\begin{equation}
\mathrm{D}\bigl(L^{2}_{\mathrm{res}}([GL_{2}])^{\mathrm{sph\text{-}fin}}\bigr)\;\simeq\;\mathrm{Sp}
\end{equation}
sends the constant function \(\mathbb{1}_{[GL_{2}]}\) to the sphere
spectrum.
\end{theorem}

\begin{proof}
The residual subspace is one-dimensional and trivial-character
(Mo{\oe}glin--Waldspurger II.2.4); hence the derived category of
its sph-fin completion is a single-object stable category
\(\mathrm{Sp}\). On the other side, the empty restricted tensor
product over all places with distinguished unit the sphere spectrum
is the sphere spectrum by the universal property of restricted
tensor in \(\mathrm{Pr}^{\mathrm{L,st}}\) (F2.c tightened). The two
match canonically.
\end{proof}

\begin{theorem}[categorical lift of \(L^{2}_{\mathrm{cont}}\) on the non-exceptional locus]
\label{v4-arith:w6-d:thm:cont-lift}
\ClaimStatusProvedHere. On the non-exceptional locus
\(\Sigma^{\mathrm{nonexc}}=i\mathbb{R}\setminus\{0\}\), the restricted
tensor product
\((\bigotimes'_{v}\mathcal{C}^{\mathrm{cont}}_{v})/\mathrm{Weyl}\)
is canonically equivalent to
\(\mathrm{D}(L^{2}_{\mathrm{cont}}([GL_{2}])^{\mathrm{sph\text{-}fin}})^{\mathrm{nonexc}}\),
the sph-fin derived category of the continuous summand restricted
to the non-exceptional unitary axis.
\end{theorem}

\begin{proof}
By \Cref{v4-arith:w6-d:thm:PW-lift}: the Plancherel synthesis of
\Cref{v4-arith:w6-d:thm:PW-synthesis} gives the Hilbert-space isomorphism
with the \(\xi\)-weighted \(L^{2}\)-space; this lifts to derived categories on
admissible / holomorphic / compactly generated classes. Restriction
to the non-exceptional locus removes the two polar points \(s=0,1\)
where \(N\) fails to be an involution; on the complement, the Weyl
quotient is a \(\mathbb{Z}/2\)-Galois quotient in
\(\mathrm{Pr}^{\mathrm{L,st}}\) and computes as a homotopy
fixed-point category (Lurie HA 4.8.5.16).
\end{proof}

\subsection{Sphericity-Finite Compatibility on the Continuum}

\begin{proposition}[sph-fin preservation under continuous-spectrum restriction]
\label{v4-arith:w6-d:prop:sph-fin-cont}
\ClaimStatusProvedHere. The sph-finiteness condition (compact
generation with finite-type local Hecke modules) is preserved
under restriction to the continuous summand and the Weyl quotient:
if every \(\mathcal{C}^{\mathrm{cont}}_{v}\) is sph-fin, then
\((\bigotimes'_{v}\mathcal{C}^{\mathrm{cont}}_{v})/\mathrm{Weyl}\)
restricted to the non-exceptional locus is sph-fin.
\end{proposition}

\begin{proof}
Sph-fin is a compact-generation + finite-type property; compact
generation is preserved under filtered colimits of compactly
generated categories (Lurie HA 5.3.2.9) and under \(\mathbb{Z}/2\)
homotopy fixed-point quotients when the action is by compact-object-
preserving functors. The normalised intertwining operator
\(N(s,\chi)\) sends a compact object at \(s\) (the structure sheaf of a
closed point) to a compact object at \(-s\), by the local holomorphic
extension (\Cref{v4-arith:w6-d:prop:local-FE}.3). Restricted to the
non-exceptional locus, \(N\) is unitary hence a compact-preserving
equivalence. Sph-fin is preserved.
\end{proof}

\section{Module (E3): Assembly Compatible with Cuspidal LACD}
\label{v4-arith:w6-d:E3}

\subsection{Full Statement}

\begin{theorem}[\(\mathrm{L}_{\mathrm{ACD}}\) on the non-exceptional
sph-finite locus, Eisenstein case]
\label{v4-arith:w6-d:thm:LACD-Eisenstein-nonexc}
\ClaimStatusProvedHereConditional. On the non-exceptional sph-finite locus,
there is a canonical equivalence
\begin{equation}
\label{v4-arith:w6-d:eq:LACD-Eisenstein-nonexc}
\Bigl(\bigotimes_{v}^{\prime}\mathcal{C}_{v}\Bigr)^{\mathrm{nonexc}}\;\xrightarrow{\;\sim\;}\;\mathrm{D}^{\mathrm{sph\text{-}fin}}([GL_{2}])^{\mathrm{nonexc}}
\end{equation}
of stable presentable \(\infty\)-categories, where:
\begin{itemize}
\item \(\mathcal{C}_{v}=\mathrm{D}^{\mathrm{sph\text{-}fin}}([GL_{2}(\mathbb{Q}_{v})])\)
decomposes as an internal direct sum
\(\mathcal{C}^{\mathrm{cusp}}_{v}\oplus\mathcal{C}^{\mathrm{res}}_{v}\oplus\mathcal{C}^{\mathrm{cont}}_{v}\)
(local Langlands decomposition; Bernstein--Zelevinsky 1977 for
\(p\)-adic; Wallach 1988 Ch.~11 at archimedean);
\item the restricted tensor product respects this decomposition
summand by summand, with the Weyl-quotient convention on the
continuous summand;
\item the non-exceptional superscript refers to the complement of
\(\Sigma^{\mathrm{exc}}=\{s=0\}\cup\{s=1\}\) on the continuous-
spectrum factor, and to the full cuspidal and residual summands
with no restriction.
\end{itemize}
under the local categorical comparison and global restricted-tensor
identification inputs of \(\mathrm{L}_{\mathrm{ACD}}\). The equivalence
is compatible with:
\begin{enumerate}
\item local Whittaker models at each finite prime
(Bernstein--Zelevinsky 1977 \S 2);
\item Bernstein-centre actions across places
(\Cref{v4-arith:lacd:thm:L2-closed});
\item archimedean \((\mathfrak{g},K_{\infty})\)-module structure at
\(v=\infty\) (\Cref{v4-arith:lacd:thm:L0-closed});
\item normalised intertwining operators
\(N_{v}(s,\chi)\) at each place (\Cref{v4-arith:w6-d:def:N-normalised});
\item Paley--Wiener synthesis
(\Cref{v4-arith:w6-d:thm:PW-synthesis}) on the continuous summand.
\end{enumerate}
\end{theorem}

\begin{proof}
\textbf{Step 1 (cuspidal).} The cuspidal summand is supplied by the cuspidal LACD comparison
(\Cref{v4-arith:w5-e:thm:LACD-cuspidal}). Nothing new is
required; the cuspidal factor contributes through the orthogonal
direct-sum on each side.
\textbf{Step 2 (residual).} The residual summand is treated in
\Cref{v4-arith:w6-d:thm:res-lift}: restricted tensor product of
\(\mathrm{Sp}\)-factors is \(\mathrm{Sp}\); on the other side the
residual space is one-dimensional.
\textbf{Step 3 (continuous non-exceptional).} The continuous
summand restricted to \(\Sigma^{\mathrm{nonexc}}\) is controlled by
\Cref{v4-arith:w6-d:thm:cont-lift}: the categorical lift uses the
Plancherel synthesis of \Cref{v4-arith:w6-d:thm:PW-synthesis} plus
the coherent Weyl action of \Cref{v4-arith:w6-d:prop:Weyl-coherent}.
\textbf{Step 4 (assembly).} The three summands assemble through
the orthogonal decomposition of \(L^{2}([GL_{2}])\)
(\Cref{v4-arith:w6-d:thm:Langlands-decomp}); on the restricted-
tensor side they assemble through the local summand decomposition
of each \(\mathcal{C}_{v}\). The restricted tensor product of an
orthogonal direct sum is the orthogonal direct sum of restricted
tensor products (Lurie HA 4.8.2: tensor product distributes over
direct sum in \(\mathrm{Pr}^{\mathrm{L,st}}\)).
\textbf{Step 5 (compatibilities).} Whittaker / Bernstein /
archimedean compatibilities are inherited from
\Cref{v4-arith:w5-e:thm:LACD-cuspidal} on the cuspidal factor and
from \Cref{v4-arith:lacd:thm:L0-closed}--\Cref{v4-arith:lacd:thm:L3-closed}
on the archimedean and cross-prime sides. The intertwining
compatibility is the construction of \(N\) in
\Cref{v4-arith:w6-d:prop:Weyl-coherent}. Paley--Wiener
compatibility is the content of \Cref{v4-arith:w6-d:thm:PW-lift}.
\end{proof}

\subsection{Corollaries and Domain}

\begin{corollary}[adelic categorical descent reduction on the non-exceptional sph-finite locus]
\label{v4-arith:w6-d:cor:full-LACD-nonexc}
\ClaimStatusProvedHereConditional. On the non-exceptional sph-finite locus,
\(\mathrm{L}_{\mathrm{ACD}}\) (\Cref{v4-arith:acd:statement}) holds
under the standing local categorical comparison and global
restricted-tensor identification inputs. The equivalence
\begin{equation}
\bigotimes_{v}^{\prime}\mathrm{D}^{\mathrm{sph\text{-}fin}}([GL_{2}(\mathbb{Q}_{v})])\;\xrightarrow{\;\sim\;}\;\mathrm{D}^{\mathrm{sph\text{-}fin}}([GL_{2}])
\end{equation}
is valid on the complement of \(\Sigma^{\mathrm{exc}}=\{s=0\}\cup\{s=1\}\)
within the Eisenstein continuous parameter, and everywhere on the
cuspidal and residual summands.
\end{corollary}

\begin{proof}
Combine the cuspidal LACD comparison cuspidal
(\Cref{v4-arith:w5-e:thm:LACD-cuspidal}) with residual
(\Cref{v4-arith:w6-d:thm:res-lift}) and continuous non-exceptional
(\Cref{v4-arith:w6-d:thm:cont-lift}) through the assembly theorem
\Cref{v4-arith:w6-d:thm:LACD-Eisenstein-nonexc}.
\end{proof}

\begin{conjecture}[exceptional unitary-crossing Ext\({}^{1}\) match]
\label{v4-arith:w6-d:conj:exceptional-unitary-crossing}
\ClaimStatusConjectured. Essential surjectivity extends from the
non-exceptional locus to the full sph-finite automorphic category
by matching the Ext\({}^{1}\)-class of the local intertwining
data at \(s=0,1\). Specifically, at \(s=0\) (spherical crossing) the
local Ext\({}^{1}\)-class of the reducibility short exact sequence
of \(I_{v}(0,\chi)\) matches the deformation-theoretic Ext\({}^{1}\)
of the Emerton--Gee stack at the corresponding unramified
Frobenius-trace-zero point; at \(s=1\) (residual crossing) the
one-dimensional residual space matches the unique closed stratum
of the Emerton--Gee stack at the cyclotomic character. The match
is a single Ext\({}^{1}\) datum per crossing, analogous to
\Cref{v4-arith:f1a3-BH:conj:Paskunas-wild-ext} but at the
automorphic (rather than local Langlands) level.
\end{conjecture}

\begin{remark}[domain of the exceptional residual]
\label{v4-arith:w6-d:rem:exceptional-domain}
The two exceptional points are measure-zero in \(i\mathbb{R}\) and
do not affect the Plancherel measure \(|\xi(2s)|^{-2}ds\). They do
affect the ind-categorical structure: at \(s=0,1\) the local
\(\mathcal{C}^{\mathrm{cont}}_{v}\) has a non-semisimple
Ext\({}^{1}\)-extension that must be matched by an
Ext\({}^{1}\)-class in the global presentable \(\infty\)-category.
It is a non-formal theorem, parallel to the F1.a.2 non-semisimple
class-match of \Cref{v4-arith:f1a3-BH:conj:Paskunas-wild-ext}.
\end{remark}

\subsection{Compatibility with the Langlands--Shahidi normalisation}

\begin{proposition}[cross-reference to the Langlands--Shahidi normalisation Koszul--Petersson correction]
\label{v4-arith:w6-d:prop:W6C-crossref}
\ClaimStatusProvedHereConditional \emph{(conditional on the Langlands--Shahidi normalisation LS
Plancherel-measure identification)}. The Langlands--Shahidi normalised operator
\(N_{v}(s,\chi)\) of \Cref{v4-arith:w6-d:def:N-normalised} enters the
the Langlands--Shahidi normalisation Koszul normalisation correction of \Cref{v4-arith:w5-c:prop:LS-absorption}
through the absolute-value specialisation
\(N^{\mathrm{LS}}(s)=\prod_{v}N_{v}(s,\chi_{\mathrm{triv}})\)
on the critical line, where \(\chi_{\mathrm{triv}}\) is the
trivial Hecke character. The conditional Plancherel-normalisation
match of the Langlands--Shahidi normalisation corresponds on the categorical side to the
unitary-axis Weyl-symmetry of the Paley--Wiener space
\(\mathrm{PW}(i\mathbb{R},V_{\chi})\) under the coherent Weyl
action \Cref{v4-arith:w6-d:prop:Weyl-coherent}. No fresh
computation is needed beyond the stated measure-identification input.
\end{proposition}

\begin{proof}
Immediate by comparison of the Langlands--Shahidi normalisation definition
\Cref{v4-arith:w5-c:def:LS-factor} with the adelic global product
of \Cref{v4-arith:w6-d:def:N-normalised} on the critical line,
using the Tate global \(\varepsilon=1\) for \(GL_{1}\).
\end{proof}

\section{Comparison and Hypotheses}
\label{v4-arith:w6-d:source-comparison}

\subsection{Independent Inputs}

\paragraph{Derivation and comparison sources.}
The derivation of the cuspidal LACD comparison cuspidal reduction used Flath 1979,
Bernstein 1988, Mo{\oe}glin--Waldspurger 1995 \S III (cuspidal
orthogonality). The present Eisenstein reduction uses Langlands 1976,
Shahidi 1990, 2010, Arthur 2013 \S 3.5, Harish-Chandra 1976,
Wallach 1992 \S 12, and Mo{\oe}glin--Waldspurger 1995 \S II, IV
(spectral theory of Eisenstein series). The categorical assembly
uses Lurie HA 4.8 and 5.1--5.3 (restricted tensor and monoidal
coherence), F2.c tightened (\Cref{v4-arith:tate:thm:F2c-main}),
and the cuspidal LACD comparison cuspidal (\Cref{v4-arith:w5-e:thm:LACD-cuspidal}).
The cuspidal and Eisenstein derivation paths are disjoint at the
chapter level of Mo{\oe}glin--Waldspurger (\S III vs \S II/IV);
no primary author appears in both spectral-decomposition roles.

\paragraph{Intertwining-operator normalisation.}
Langlands 1976 defines the unnormalised operator \(M(s,\chi)\) by
an integral over \(N(\mathbb{A})\). Shahidi 1990, 2010 defines the
normalised operator \(N(s,\chi)\) with \(L\)- and \(\varepsilon\)-factor
corrections. Mo{\oe}glin--Waldspurger 1995 uses the Arthur
normalisation \(\tilde M(s,\chi)\) with a specific \(L\)-ratio.
The three normalisations agree on the critical line after the
Tate global product \(\prod_v \varepsilon_v = 1\) on \(GL_1\); the
Weyl translation \(\chi\mapsto\chi^{w_0}\) is consistent across
all three.

\paragraph{Categorical ambient structure.}
Each \(\mathcal{C}^{\mathrm{cont}}_{v}\) is a stable presentable
\(\infty\)-category (\Cref{v4-arith:w6-d:prop:C-cont-stable-pres});
restricted tensor products lie in \(\mathrm{Pr}^{\mathrm{L,st}}\)
(F2.c tightened); Weyl quotients are homotopy fixed-point
constructions in \(\mathrm{Pr}^{\mathrm{L,st}}\) (Lurie HA 4.8.5.16).
The Langlands spectral decomposition
\Cref{v4-arith:w6-d:thm:Langlands-decomp} is an orthogonal direct
sum at both Hilbert-space and derived-category level.

\paragraph{Hypotheses.}
The non-exceptional hypothesis is made explicit in
\Cref{v4-arith:w6-d:def:exceptional}. Hecke-character unitarity
is required for unitarity of \(N_v\)
(\Cref{v4-arith:w6-d:prop:local-FE}.2); it is imposed throughout.
The spherical / Iwahori-finite distinguished-unit hypothesis carries
over from the cuspidal LACD comparison through the construction of \(u^{\mathrm{cont}}_v\).

\paragraph{Functoriality of the Weyl action.}
The intertwining operators act place by place; cross-place
commutation is \Cref{v4-arith:w6-d:prop:Weyl-coherent}. The Weyl
action is a coherent \(\mathbb{Z}/2\) action on the restricted tensor
product, not an \(\infty\)-groupoid action; coherence up to homotopy
follows from Lurie HA 5.1.2.26.

\paragraph{Open input at the exceptional locus.}
The exceptional unitary-crossing \(\mathrm{Ext}^1\)-class match is
\Cref{v4-arith:w6-d:conj:exceptional-unitary-crossing}; it is named
as a conjecture. On the non-exceptional locus, every input is
either a primary theorem (the Plancherel synthesis of
\Cref{v4-arith:w6-d:thm:PW-synthesis}, sourced from Shahidi 2010
Thm~4.2.2) or a proved claim in this chapter.

\paragraph{Numerical comparison.}
This chapter contains no numerical computation; the expected-value
comparison is through the analytic Plancherel formula.

\subsection{Composite with the cuspidal LACD comparison}

\begin{theorem}[composite \(\mathrm{L}_{\mathrm{ACD}}\) reduction after the cuspidal LACD comparison + the Eisenstein LACD comparison]
\label{v4-arith:w6-d:thm:composite-LACD}
\ClaimStatusProvedHereConditional. Combining the cuspidal LACD comparison cuspidal reduction
\Cref{v4-arith:w5-e:thm:LACD-cuspidal} with the Eisenstein LACD comparison residual +
continuous non-exceptional reduction
\Cref{v4-arith:w6-d:thm:LACD-Eisenstein-nonexc},
\(\mathrm{L}_{\mathrm{ACD}}\) is reduced on the sph-finite automorphic
category modulo the two unitary-axis crossings
\(\Sigma^{\mathrm{exc}}=\{s=0\}\cup\{s=1\}\), and modulo the standing
local categorical comparison, global restricted-tensor identification,
and BC-diamond-glue inputs.
\end{theorem}

\begin{proof}
Orthogonal direct-sum assembly of the three reductions through
\Cref{v4-arith:w6-d:thm:Langlands-decomp} and the restricted-tensor
orthogonality of Lurie HA 4.8.
\end{proof}

\section{F1 Obstruction Decomposition on the Eisenstein Continuum}
\label{v4-arith:w6-d:residual-table}

\begin{proposition}[F1 obstruction decomposition after the Eisenstein LACD comparison]
\label{v4-arith:w6-d:prop:F1-table-after-w6d}
\ClaimStatusProvedHere. The Eisenstein LACD comparison and the
cuspidal obstruction decomposition
\Cref{v4-arith:w5-e:prop:F1-table-refined} give the
F1 obstruction decomposition is:
\begin{enumerate}
\item \textbf{F1.a.1.non-ord} (non-formal theorem): essential surjectivity
of DEG at \(p\geq 5\) on the non-ordinary locus.
\item \textbf{F1.a.1.TW-removal} (non-formal theorem): extension beyond
Taylor--Wiles hypotheses.
\item \textbf{F1.a.2.class-match} (non-formal theorem): class-level
Pa{\v s}k{\=u}nas Ext\({}^{1}\)-compatibility at \(p=2,3\).
\item \textbf{\(\mathrm{L}_{\mathrm{ACD}}\).exceptional-crossing}
(non-formal theorem): Ext\({}^{1}\)-class match at \(s=0,1\)
(\Cref{v4-arith:w6-d:conj:exceptional-unitary-crossing}); the Eisenstein LACD comparison
replaces the non-formal theorem Eisenstein conjecture
\Cref{v4-arith:w5-e:conj:LACD-Eisenstein} with this narrower
residual.
\item \textbf{\(\mathrm{L}_{\mathrm{ACD}}\).local-global comparison}:
the local categorical comparison and the global restricted-tensor
identification of \Cref{v4-arith:acd:residuals}.
\item \textbf{BC-diamond-glue} (non-formal theorem):
Bernstein-centre diamond-gluing, Chapter~\ref{v4-arith:bcglue-core}.
\end{enumerate}
No numerical closure total is asserted here.
\end{proposition}

\begin{proof}
Tabulation. the cuspidal LACD comparison carried the Eisenstein residual at non-formal theorem;
the Eisenstein LACD comparison reduces the non-exceptional part and isolates a non-formal theorem
Ext\({}^{1}\)-class match at the two polar crossings. The global LACD
comparison inputs remain.
\end{proof}

\section{Eisenstein Boundary}
\label{v4-arith:w6-d:summary}

\begin{enumerate}
\item \textbf{Langlands--Shahidi normalisation constructed}
(\Cref{v4-arith:w6-d:def:N-normalised},
\Cref{v4-arith:w6-d:prop:local-FE}): normalised intertwining
operators \(N_{v}(s,\chi)\) with functional equation, unitarity on
the axis, polar divisor at \(s=0,1\). Compatible with the Langlands--Shahidi normalisation through
\Cref{v4-arith:w6-d:prop:W6C-crossref} and the conditional normalisation match of
\Cref{v4-arith:w5-c:prop:LS-absorption}.
\item \textbf{Paley--Wiener synthesis categorically lifted}
(\Cref{v4-arith:w6-d:thm:PW-lift}): the continuous summand of
\(L^{2}([GL_{2}])\) lifts to a restricted tensor product of local
continuous-spectrum presentable stable \(\infty\)-categories modulo
a coherent Weyl quotient.
\item \textbf{Langlands spectral decomposition categorified}
(\Cref{v4-arith:w6-d:thm:Langlands-decomp},
\Cref{v4-arith:w6-d:thm:res-lift},
\Cref{v4-arith:w6-d:thm:cont-lift}): three-summand orthogonal
direct sum at Hilbert-space and derived-category level, with
residual summand = \(\mathrm{Sp}\), continuous summand =
Paley--Wiener sheaves on the unitary axis.
\item \textbf{\(\mathrm{L}_{\mathrm{ACD}}\) reduced on the
non-exceptional sph-finite locus}
(\Cref{v4-arith:w6-d:thm:LACD-Eisenstein-nonexc},
\Cref{v4-arith:w6-d:cor:full-LACD-nonexc}): the restricted tensor
product equals the global sph-finite automorphic derived category
on the complement of two unitary-axis crossings.
\item \textbf{Exceptional residual isolated}
(\Cref{v4-arith:w6-d:conj:exceptional-unitary-crossing}):
Ext\({}^{1}\)-class match at \(s=0\) (spherical crossing) and \(s=1\)
(residual crossing); non-formal theorem of domain, parallel to the
Pa{\v s}k{\=u}nas non-semisimple class-match of F1.a.2.
\item \textbf{Composite reduction}
(\Cref{v4-arith:w6-d:thm:composite-LACD}): the cuspidal LACD comparison cuspidal +
the Eisenstein LACD comparison residual and non-exceptional continuous jointly reduce
\(\mathrm{L}_{\mathrm{ACD}}\) on the sph-finite automorphic category
modulo the two-point exceptional residual and the standing comparison
inputs.
\item \textbf{Obstruction decomposition}
(\Cref{v4-arith:w6-d:prop:F1-table-after-w6d}): F1 residuals drop
to the named items above.
\end{enumerate}

The Eisenstein / continuous-spectrum residual
\Cref{v4-arith:w5-e:conj:LACD-Eisenstein} is reduced on its
natural domain (non-exceptional sph-finite); the two remaining
unitary-axis crossings are isolated as the narrow
\(\mathrm{Ext}^1\)-class-match conjecture
\Cref{v4-arith:w6-d:conj:exceptional-unitary-crossing}, a 10--15
page condition. the Langlands--Shahidi normalisation states the Plancherel-normalisation match from the
Koszul-pairing side; this chapter establishes the categorical lift
from the intertwining-operator side; their cross-compatibility is
\Cref{v4-arith:w6-d:prop:W6C-crossref}.

The Eisenstein open condition is
\Cref{v4-arith:w6-d:conj:exceptional-unitary-crossing}: matching
\(\mathrm{Ext}^1\)-classes at \(s=0\) and \(s=1\), where the
reducibility of \(I_v(s,\chi)\) and the Emerton--Gee stack geometry
must be reconciled. The global LACD comparison table remains separate.
